\section{Detection Sensitivity and the Bayesian Update}
\label{sec:detection_sensitivity}

The Bayesian update in Section~\ref{sec:relaxation_protocol}
depends on the detector's false-negative rate $\beta(\testdur)$
over the test window.  This section characterizes $\beta$ in
terms of the multi-channel detection architecture
of~\citep[Table~1]{lasser2026scm} and records the per-test
evidence strength that follows.

\subsection{The detection model}

Within a single test, $A_n = 1$ means the pathway activates at
some point during the $\testdur$-step window; $O_n = 1$ means the
detector records contraction during that window.  The detection
sensitivity $\beta$ mediates between them:
$\beta = \Pr(O_n = 0 \mid A_n = 1)$.

\begin{assumption}[Activation Duration]
\label{ass:D0}
When the pathway activates during a test ($A_n = 1$), the
activation persists for $k_{\mathrm{act}}$ consecutive steps,
where $k_{\mathrm{act}}$ is either (a)~a fixed parameter of the
pathway (determined by
~\citep[Definition~2]{lasser2026horizon}'s exercise-pathway
structure), or (b)~drawn from a distribution $K$ with known or
bounded support.  The per-test false-negative rate $\beta$
depends on $k_{\mathrm{act}}$: the detector has
$\min(k_{\mathrm{act}}, \testdur)$ steps to detect activation.
For the default analysis we assume
$k_{\mathrm{act}} \geq \testdur$ (persistent activation); the
transient case $k_{\mathrm{act}} < \testdur$ is handled as a
sensitivity analysis below.
\end{assumption}

\begin{definition}[Detection Sensitivity]
\label{def:detection_sensitivity}
For a test $T = (R, \Xirestrict, \scope, \testdur, \damagebound)$
monitoring pathway $\Pathway = (\pi_1, \ldots, \pi_L)$, the
per-step detection probability on channel $c_i$ is
$$
\delta_i
   = \Pr(c_i \text{ detects activation at } t'
         \mid \Pathway \text{ activates at } t').
$$
The aggregate per-step detection probability across all active
channels is
\begin{equation}
\label{eq:delta_agg}
\deltaagg = 1 - \prod_{c_i \in \mathcal{O}(P_t) \cap \mathcal{C}_{\scope}}
  (1 - \delta_i),
\end{equation}
where $\mathcal{C}_{\scope}$ is the set of channels that cover
the scope $\scope$.
\end{definition}

\begin{assumption}[Cross-Channel Independence]
\label{ass:D1}
Detection events on distinct channels are conditionally
independent given the latent activation state $A_n$.  This is
justified for cross-substrate channels by channel
isolation~\citep[Proposition~5]{lasser2026exo} (unconditional),
and for same-substrate channels under the joint-factorization
assumption of~\citep[Proposition~6]{lasser2026exo} (conditional
on committed verification).  When
Assumption~\ref{ass:D1} fails (correlated same-substrate
channels without committed verification), the product
formula~\eqref{eq:delta_agg} is optimistic; a conservative
fallback uses only the cross-substrate channels,
\begin{equation}
\label{eq:delta_agg_conservative}
\deltaagg^{\mathrm{cons}}
  = 1 - \prod_{c_i \in \mathcal{O}(P_t) \cap \mathcal{C}_{\scope}
              :\ \mathrm{cross\text{-}substrate}} (1 - \delta_i).
\end{equation}
\end{assumption}

\begin{assumption}[Temporal Independence or Mixing]
\label{ass:D2}
Conditional on the latent activation state $A_n$, detection
outcomes at distinct time steps within a test window are either
(a)~independent (memoryless detectors) or (b)~satisfy a mixing
condition with mixing time
$\testmix \ll \testdur$.  Without Assumption~\ref{ass:D2},
persistent detector outages or common-cause failures can produce
correlated misses across time steps, making $\beta$ decay slower
than exponential.  Under (b), the effective number of
independent detection opportunities in $\testdur$ steps is
$\testdur / \testmix$ rather than $\testdur$.
\end{assumption}

Under Assumptions~\ref{ass:D0}--\ref{ass:D2} with persistent
activation, the window-level false-negative rate satisfies
\begin{equation}
\label{eq:beta_window}
\beta(\testdur) \leq (1 - \deltaagg)^{\testdur / \testmix}.
\end{equation}
Under the memoryless case (\ref{ass:D2}(a)), $\testmix = 1$ and
$\beta(\testdur) = (1 - \deltaagg)^{\testdur}$.  For transient
pathways with activation duration $k < \testdur$, replace
$\testdur$ by $\min(k, \testdur)$ throughout.

\subsection{Detection-sensitivity bounds}

\begin{proposition}[Detection Sensitivity Bounds]
\label{prop:detection_bounds}
Under the multi-channel architecture
of~\citep[Table~1]{lasser2026scm} and
Assumptions~\ref{ass:D0}--\ref{ass:D2}:
\begin{itemize}
\item[(a)] $\beta(\testdur)$ is monotonically decreasing in
  $\testdur$: longer tests have lower false-negative rates.
\item[(b)] For any $\epsilon_\beta \in (0, 1)$, the minimum test
  duration required to achieve $\beta(\testdur) \leq
  \epsilon_\beta$ is
  \begin{equation}
  \label{eq:tau_min}
  \testmin = \left\lceil
    \frac{\log \epsilon_\beta}{\log(1 - \deltaagg)}
  \right\rceil,
  \end{equation}
  rescaled by $\testmix$ under the mixing case of
  Assumption~\ref{ass:D2}(b).
\item[(c)] If at least one cross-substrate channel covers
  $\scope$ with detection probability $\deltacross > 0$, then
  even without Assumption~\ref{ass:D1} for same-substrate
  channels,
  \begin{equation}
  \label{eq:cross_substrate_bound}
  \beta(\testdur) \leq (1 - \deltacross)^{\testdur / \testmix},
  \end{equation}
  which converges to zero exponentially in $\testdur$ at rate
  $1/\testmix$.
\end{itemize}
\end{proposition}

\begin{proof}[Proof sketch]
Part~(a) is immediate from~\eqref{eq:beta_window}: adding more
independent detection opportunities monotonically decreases the
product.  Part~(b) inverts the bound
in~\eqref{eq:beta_window}.  Part~(c) uses that cross-substrate
channels satisfy Assumption~\ref{ass:D1} unconditionally
by~\citep[Proposition~5]{lasser2026exo}, so the product formula
is exact for them; temporal independence follows from
Assumption~\ref{ass:D2}(a) or (b).
\end{proof}

Part~(c) is the substrate-diversity dividend for detection:
cross-substrate channels provide independent detection
opportunities that same-substrate channels may not (due to
correlated failures), and their isolation is unconditional.

\begin{remark}[Cross-substrate-only rate recovery]
\label{rem:cross_substrate_rate}
The interaction between Assumptions~\ref{ass:D1}
(cross-channel independence) and~\ref{ass:D2} (temporal
independence) is subtle.  A common-cause failure mode that
violates~\ref{ass:D1} for same-substrate channels (e.g., a shared
software dependency that miscalibrates an entire silicon
detector cohort) typically also violates~\ref{ass:D2} for the
same-substrate temporal stream (the failure persists across
time).  The actually-dangerous regime is therefore
\emph{joint} same-substrate failure across channels and across
time.  The conservative bound $\deltaagg^{\mathrm{cons}}$
of~\eqref{eq:delta_agg_conservative} protects against both
simultaneously: cross-substrate channels, by~\citep[Proposition~5]{lasser2026exo}'s
unconditional channel-isolation, satisfy~\ref{ass:D1}
\emph{and}~\ref{ass:D2}(a) without further assumption, because
the temporal failure of a biological observer is independent of
the temporal failure of a silicon observer.  Under
$\deltaagg^{\mathrm{cons}} > 0$, Theorem~\ref{thm:convergence}'s
KL rate is recovered up to a factor that depends only on the
ratio of cross-substrate to all-channel detection probability:
$\KL^{\mathrm{cons}} / \KL^{\mathrm{full}} \geq
\deltaagg^{\mathrm{cons}} / \deltaagg$ in the regime
$\deltaagg \ll 1$, with equality in the limit of a single
dominant channel.  Convergence is preserved unconditionally; the
constant factor degradation is the cost of the joint failure
mode.
\end{remark}

\subsection{Composition rule for transient pathways}
\label{sec:transient_composition}

Assumption~\ref{ass:D0} treats persistent activation as the
default ($k_{\mathrm{act}} \geq \testdur$).  For transient
pathways with activation duration $k < \testdur$, the
window false-negative rate is
$\beta(\testdur) = (1 - \deltaagg)^{\min(k, \testdur)}$ under
Assumption~\ref{ass:D2}(a), or scaled by $\testmix$ under
Assumption~\ref{ass:D2}(b).  When $k = 1$ (single-step
activation) the formula gives $\beta = 1 - \deltaagg$, which
may be non-negligible.  An operational composition rule fixes
the design parameter.

\begin{proposition}[Detection design rule for transient pathways]
\label{prop:transient_design}
For a target window false-negative rate $\epsilon_\beta$ and a
known (or upper-bounded) pathway duration $k$:
\begin{itemize}
\item Under Assumption~\ref{ass:D2}(a) (memoryless detectors),
  the per-step detection probability $\deltaagg$ during the
  test must satisfy
  $(1 - \deltaagg)^k \leq \epsilon_\beta$, i.e.,
  $\deltaagg \geq 1 - \epsilon_\beta^{1/k}$.  For $k = 1$ this
  collapses to $\deltaagg \geq 1 - \epsilon_\beta$ (the
  single-step detector must itself meet the target); for larger
  $k$ the requirement softens proportionally.
\item Under Assumption~\ref{ass:D2}(b) (mixing with mixing
  time $\testmix$), the effective number of independent
  detection opportunities in $k$ steps is $k/\testmix$, and the
  rule strengthens to $\deltaagg \geq 1 -
  \epsilon_\beta^{\testmix/k}$; recover memoryless as
  $\testmix \to 1$.
\end{itemize}
Implementations that cannot raise $\deltaagg$ at fixed $\testdur$
must either extend the test window (use the persistent-activation
form with longer $\testdur$) or raise the detector sampling rate
within the window (effectively shrinking $\testmix$).  The choice
between these is implementation-defined; the composition rule
sets the constraint each must satisfy.
\end{proposition}

\subsection{Evidence strength of a single test}

\begin{proposition}[Evidence Strength of Controlled Relaxation]
\label{prop:evidence_strength}
After a controlled relaxation test $T$ with outcome
$O_n \in \{0, 1\}$, using the binary model of
Section~\ref{sec:relaxation_protocol} ($H_0$: $\theta = 0$ vs.\
$H_1$: $\theta = \thetaone$ with fixed likelihoods):
\begin{itemize}
\item[(a)] \textbf{Evidence for $R$ ($O_n = 1$, contraction
  detected).}
  \begin{equation}
  \label{eq:log_L1}
  \Delta \logit = \log \Lone
    = \log \frac{\thetaone (1 - \beta) + (1 - \thetaone) \alphafp}{\alphafp}.
  \end{equation}
  This is a positive constant under
  Assumption~\ref{ass:T1} ($\Lone > 1$).  For the common case
  $\thetaone = 1$: $\Delta \logit = \log((1 - \beta)/\alphafp)$.
\item[(b)] \textbf{Evidence against $R$ ($O_n = 0$, no
  contraction detected).}
  \begin{equation}
  \label{eq:log_L0}
  \Delta \logit = \log \Lzero
    = \log \frac{\thetaone \beta + (1 - \thetaone)(1 - \alphafp)}{1 - \alphafp}.
  \end{equation}
  This is a negative constant under
  Assumption~\ref{ass:T1} ($\Lzero < 1$).  For $\thetaone = 1$:
  $\Delta \logit = \log(\beta / (1 - \alphafp))$.  The
  denominator $(1 - \alphafp)$ is material when $\alphafp$ is
  non-negligible and was missing in a naive derivation that
  assumed $\alphafp = 0$.
\item[(c)] \textbf{Comparison with the restricted regime.}
  Under the restricted regime
  (Proposition~\ref{prop:asym_evidence}), the capability is not
  exercised, so both outcomes have likelihood ratio~$1$ (no
  update).  Under controlled relaxation the capability is
  exercised
  (Assumption~\ref{ass:T5}), so both $\Lzero$ and $\Lone$ are
  non-unit and both outcomes are informative.
\item[(d)] \textbf{Effect of $\thetaone$.}
  The alternative hypothesis $\thetaone$ affects evidence
  strength:
  \begin{itemize}
  \item $\thetaone = 1$: $\Lone = (1-\beta)/\alphafp$, $\Lzero =
    \beta/(1-\alphafp)$.  Maximum evidence per test.
  \item $\thetaone < 1$: $\Lone$ decreases and $|\log \Lzero|$
    decreases.  Less evidence per test, but the alternative may
    be more realistic for stochastically activating pathways.
  \end{itemize}
  The choice of $\thetaone$ is a modelling decision that does
  not affect convergence (Theorem~\ref{thm:convergence} holds
  for any $\thetaone \in (0, 1]$ under
  Assumption~\ref{ass:T1}), only the rate.
\end{itemize}
\end{proposition}

\begin{proof}[Proof sketch]
Direct from the fixed likelihood ratios~\eqref{eq:L_defs}.  The
log-odds form follows from~\eqref{eq:update_pos} and
\eqref{eq:update_neg}: $\logit \pR(t^+) = \logit \pR(t) +
\log(\Pr(O_n \mid H_1) / \Pr(O_n \mid H_0))$.
\end{proof}

\begin{remark}[Integration with log-odds aggregation]
\label{rem:log_odds_integration}
The log-odds evidence from
Proposition~\ref{prop:evidence_strength} feeds directly into the
log-odds risk aggregation
of~\citep[Definition~5]{lasser2026scm} as an additional
evidence channel.  The controlled-relaxation test produces
evidence with the same structure as an individual agent's risk
assessment: a log-odds contribution weighted by trust.  The
trust weight on this evidence is~$1$ (the observation is direct,
not mediated by an agent's report), making test outcomes the
highest-trust evidence the aggregation can receive.  This is by
design: direct observation under controlled conditions is the
gold standard for empirical risk assessment.
\end{remark}
